%%%%%解析几何模版%%%%%
\chapter*{解析几何每日一题}

\refer{《高等数学》同济版}

\exercise{一向量的终点在点$B\qty(2,-1,7)$，它在$x$轴，$y$轴，$z$轴上的投影依次为
$4,\ -4,\ 7$求这向量的起点$A$的坐标.}

\teacher{马心怡}

\postdate{2023年4月15日}

\result{\qty(-2,3,0)}

\analysis{此题考查的是空间向量运算的坐标表示，因为$\vv{a}=a_x\boldsymbol{i}+a_y\boldsymbol{j}+a_k\boldsymbol{k}$中$a_x,\ a_y,\ a_z$分别在$x,y,z$轴上投影，由此可解此题.}

\solution{
    设向量的起点$A$的坐标为$\qty(x,y,z)$.由
    \formula{\qty(2-x,-1-y,7-z)=\qty(4,-4,7)}
    可得
    \formula{
        \left\{
        \addtolength{\jot}{-1.5em}
        \begin{aligned}
            &2-x=4\\
            &-1-y=-4\\
            &7-x=7
        \end{aligned}
        \right.
    }
    求得
    \formula{x=-2,\ y=3,\ z=0}
    即起点$A$的坐标为$\qty(-2,3,0)$
}